% Part: first-order-logic
% Chapter: natural-deduction
% Section: derivations

\documentclass[../../../include/open-logic-section]{subfiles}

\begin{document}

\iftag{FOL}
      {\olfileid{fol}{ntd}{der}}
      {\olfileid{pl}{ntd}{der}}

\olsection{\usetoken{P}{derivation}}

\begin{explain}
పరికల్పన అంటే ఏమిటో చెప్పాం; నిగమన నియమాలనూ ఇచ్చాం. సహజ నిగమనంలోని
\tetoken{వ్యుత్పత్తులు}{!!^{derivation}s} వీటి నుంచి ఆగమనాత్మకంగా ఉత్పన్నమవుతాయి: ప్రతి
\tetoken{వ్యుత్పత్తి}{!!{derivation}} తనంతట తానే ఒక పరికల్పన అయినా అవుతుంది, లేదా ఒకటి,
రెండు, మూడు \tetoken{వ్యుత్పత్తులు}{!!{derivation}s} తర్వాత ఒక సరైన నిగమనం గలదైనా అవుతుంది.
\end{explain}

\begin{defn}[\tetoken{వ్యుత్పత్తి}{!!^{derivation}}]
పరికల్పనలు~$\Gamma$ నుంచి \tetoken{ఒక వాక్యం}{!!a{sentence}}~$!A$ యొక్క
\emph{\Article{derivation} \tetoken{వ్యుత్పత్తి}{!!{derivation}}} అంటే కింది షరతులను తీర్చే
\tetoken{వాక్యాలు}{!!{sentence}s} పరిమిత వృక్షం:
\begin{enumerate}
\item వృక్షంలో అన్నిటికంటే పైనున్న \tetoken{వాక్యాలు}{!!{sentence}s}, $\Gamma$లో ఉండాలి
  లేదా వృక్షంలోని ఒక నిగమనం ద్వారా \tetoken{ఉపసంహరించిన}{!!{discharged}} అయి ఉండాలి.
\item వృక్షంలో అన్నిటికంటే కిందనున్న \tetoken{వాక్యం}{!!{sentence}}~$!A$.
\item అడుగునున్న $!A$ తప్ప వృక్షంలోని ప్రతి \tetoken{వాక్యం}{!!{sentence}}, ఆ
  \tetoken{వాక్యానికి}{!!{sentence}} నేరుగా
  కింద నిలిచే నిష్కర్ష గల నిగమన నియమాన్ని సరిగ్గా వర్తింపజేయడంలో ఒక
  పూర్వాధారం.
\end{enumerate}
అప్పుడు $!A$ను ఆ \tetoken{వ్యుత్పత్తి}{!!{derivation}} యొక్క \emph{నిష్కర్ష} అనీ,
$\Gamma$ను దాని \tetoken{ఉపసంహరించని}{!!{undischarged}} పరికల్పనలు అనీ అంటాం.

$\Gamma$ నుంచి $!A$ యొక్క ఒక \tetoken{వ్యుత్పత్తి}{!!a{derivation}} ఉంటే, $\Gamma$ నుంచి
$!A$ \emph{\tetoken{వ్యుత్పాదించదగిన}{!!{derivable}}} అని అంటాం; సంకేతాల్లో: $\Gamma \Proves
!A$. ప్రతి పరికల్పనా \tetoken{ఉపసంహరించిన}{!!{discharged}} అయిన $!A$ యొక్క ఒక \tetoken{వ్యుత్పత్తి}{!!a{derivation}}
ఉంటే $\Proves !A$ అని రాస్తాం.
\end{defn}

\begin{ex}
ప్రతి పరికల్పనా తనంతట తానే ఒక \tetoken{వ్యుత్పత్తి}{!!a{derivation}}. కాబట్టి, ఉదాహరణకు,
$!A$ తనంతట తానే ఒక \tetoken{వ్యుత్పత్తి}{!!a{derivation}}; అలాగే $!B$ కూడా. వీటికి
$\Intro{\land}$ నియమాన్ని వర్తింపజేసి కొత్త \tetoken{వ్యుత్పత్తిని}{!!{derivation}}
పొందవచ్చు:
\begin{prooftree}
\AxiomC{$!A$}
\AxiomC{$!B$}
\RightLabel{\Intro{\land}}
\BinaryInfC{$!A \land !B$}
\end{prooftree}
ఈ నియమాలు సాధారణమైనవి: అందులోని $!A$, $!B$ల స్థానంలో ఏ
\tetoken{వాక్యాలనైనా}{!!{sentence}s}, ఉదాహరణకు $!C$, $!D$లను పెట్టవచ్చు. అప్పుడు
నిష్కర్ష $!C \land !D$ అవుతుంది; కాబట్టి
\begin{prooftree}
\AxiomC{$!C$}
\AxiomC{$!D$}
\RightLabel{\Intro{\land}}
\BinaryInfC{$!C \land !D$}
\end{prooftree}
ఒక సరైన \tetoken{వ్యుత్పత్తి}{!!{derivation}}. పరికల్పనల స్థానాలను కూడా మార్చవచ్చు; అప్పుడు
$!D$ అనేది $!A$ పాత్రను, $!C$ అనేది $!B$ పాత్రను పోషిస్తాయి. కనుక
\begin{prooftree}
\AxiomC{$!D$}
\AxiomC{$!C$}
\RightLabel{\Intro{\land}}
\BinaryInfC{$!D \land !C$}
\end{prooftree}
కూడా ఒక సరైన \tetoken{వ్యుత్పత్తి}{!!{derivation}}.

ఇప్పుడు మరొక నియమాన్ని, ఉదాహరణకు $\Intro{\lif}$ను, వర్తింపజేయవచ్చు.
ఇది ఒక షరతీయాన్ని నిష్కర్షగా పొందనిస్తుంది; ఆ షరతీయపు పూర్వాంగంతో
తాదాత్మ్యం గల ఏ పరికల్పననైనా \tetoken{ఉపసంహరణ}{!!{discharge}} చేయనిస్తుంది. కాబట్టి
కిందివి రెండూ సరైన \tetoken{వ్యుత్పత్తులు}{!!{derivation}s}:
\begin{prooftree}
\AxiomC{$\Discharge{!C}{1}$}
\AxiomC{$!D$}
\RightLabel{\Intro{\land}}
\BinaryInfC{$!C \land !D$}
\DischargeRule{\Intro{\lif}}{1}
\UnaryInfC{$!C \lif (!C \land !D)$}
\DisplayProof\bottomAlignProof
\AxiomC{$!C$}
\AxiomC{$\Discharge{!D}{1}$}
\RightLabel{\Intro{\land}}
\BinaryInfC{$!C \land !D$}
\DischargeRule{\Intro{\lif}}{1}
\UnaryInfC{$!D \lif (!C \land !D)$}
\end{prooftree}
అవి వరుసగా $!D \Proves !C \lif (!C \land !D)$ అనీ,
$!C \Proves !D \lif (!C \land !D)$ అనీ చూపుతాయి.

పరికల్పనలను ఉపసంహరించడం ఒక అనుమతి, తప్పనిసరి షరతు కాదని గుర్తుంచుకోండి:
పరికల్పనలను తప్పనిసరిగా ఉపసంహరించనవసరం లేదు. ముఖ్యంగా,
\tetoken{వ్యుత్పత్తిలో}{!!{derivation}} ఆ పరికల్పనలు లేకపోయినా ఒక నియమాన్ని వర్తింపజేయవచ్చు.
ఉదాహరణకు, \tetoken{ఉపసంహరించిన}{!!{discharged}} చేయడానికి $!A$ పరికల్పన లేకపోయినా కిందిది
అనుమతించబడుతుంది:
\begin{prooftree}
  \AxiomC{$!B$}
  \DischargeRule{\Intro{\lif}}{1}
  \UnaryInfC{$!A \lif !B$}
\end{prooftree}
\end{ex}

\end{document}
